\section{Discussion}
\label{sec:discussion}

The central result is architectural: for this class of recurrent attention agent, coupling between the policy and value pathways is a structural prerequisite for learning, on par with the choice of learning signal itself. A reasonable-looking modular design --- two specialised recurrent streams cleanly separated for policy and value --- does not merely underperform when the streams are isolated; it fails to acquire the task entirely. The margin between the coupled and uncoupled designs is among the largest one can produce by a single structural edit: a $\sim$$79$-point hit-rate gap and a collapse from a calibrated detector (accuracy $0.868$, hit rate $0.794$) to an agent that never once declares a change ($0.000$ hits across $500$ episodes). Two jointly necessary conditions bound the effect: the streams must share a substrate, and the image-grounded priority stream must command the decision.

\subsection{A reinforcement-learning reading: credit assignment closes or opens the loop}
The categorical nature of the failures is most naturally read through the actor--critic credit-assignment problem. In the coupled configuration the shared deep memory forces the critic to evaluate the \emph{same} evolving state the actor acts from, and the actor's own output writes that state. ``Pressing now looks good'' is then a statement about the state the actor actually occupies; credit lands where it was earned, and the actor can discover that declaring at change onset is rewarded. When the coupling is severed, the actor acts from a private memory while the critic values a different private memory; the two representations drift apart, and the advantage handed to the actor is computed on the wrong state. Credit is systematically misattributed, and the policy retreats to the only basin that is safe under any value estimate --- never act, never be penalised --- and stays there. The always-wait collapse is the rational endpoint of an agent that cannot tell when acting helps; it is a learning failure, not a capacity failure, since the network retains the parameters and pathways to express the task. The swapped-routing control fails for a complementary reason: value and policy are coupled, but the actor is now fed the value stream, whose residual is memory rather than image and which therefore carries almost no live visual content. However clean the credit assignment, there is nothing in the actor's input that can see the change, so there is nothing to which credit can be assigned. The two controls thus dissociate the two necessary conditions cleanly.

\subsection{Mapping onto the primate attention and value systems}
The working loop has a natural counterpart in the biology. A grounded priority and commitment system --- superior colliculus together with the visual cells of the frontal eye field --- commits the action and is causally necessary for covert selection: SC inactivation produces a large covert-attention deficit while cortical gain is left intact \citep{zenon2012attention, lovejoy2010inactivation}, is more devastating than FEF inactivation \citep{bollimunta2018fefsc}, and acts as the threshold that converts cortical evidence accumulation into a committed response \citep{stine2023differential}; caudate attentional modulation itself depends on SC, implicating a colliculo-striatal--cortical loop \citep{herman2020attention}. During covert attention FEF carries the signal on its image-grounded visual cells while movement cells are suppressed \citep{thompson2005neuronal, moore2001control, moore2004microstimulation}, and the decision requires grounded sensory evidence accumulated to a bound \citep{gold2007neural}. The single closest parallel to the swapped-routing failure is the finding that inactivating sensory cortex impairs discrimination while inactivating the downstream accumulator does not \citep{katz2016dissociated}: remove the grounded representation and the decision fails, exactly as when the actor is fed the memory stream. The value side maps onto the actor--critic organisation of the basal ganglia \citep{odoherty2004dissociable, maia2009reinforcement}, with the loop closed by a dopaminergic reward-prediction-error teaching signal from SNc \citep{schultz1997neural, hollerman1998dopamine} acting through a corticostriatal three-factor plasticity rule \citep{morita2014striatal}. That value and decision signals can ride a single substrate is itself observed in LIP \citep{platt1999neural, sugrue2004matching}, supporting the shared-memory premise. The superior colliculus is best understood here as a selective-attention and commitment node rather than an evaluator \citep{krauzlis2018selective}.

\subsection{Honesty about the mapping}
The correspondence is conceptual motivation, not a literal isomorphism, and its limits must be stated plainly. First, and most importantly, the substrate and direction of the critic-to-actor coupling differ between model and brain. In the model the coupling is a \emph{shared recurrent working memory} that the actor writes and both streams read; in the brain the value-to-actor influence is a \emph{broadcast dopamine scalar} flowing from the critic system to the actor, not a shared writable state. Both a write direction (priority $\rightarrow$ shared memory) and a teaching direction (critic gradient $\rightarrow$ actor) exist in the model, so it is genuinely a closed loop; but the brain analogue should be read as the dopaminergic prediction-error precedent for ``learning requires the loop,'' not as an architectural match. Second, the dorsal/ventral actor--critic split is itself contested --- action-value signals appear broadly across striatum --- so we present it as an organising idealisation rather than settled anatomy \citep{maia2009reinforcement}. Third, the requirement that the grounded stream command the action is a moderate and partly contested parallel: reward and prior information can in some studies enter the sensory or drift-rate pathway directly, so the dissociation is not absolute. Fourth, the priority and value streams do \emph{not} map onto two distinct attention nodes: the priority stream resembles the FEF-visual and SC commitment integrator, while the value stream corresponds to the value/critic system of the basal ganglia and dopamine, which is not an attention node. LIP, FEF and SC are each a single priority map, not an actor/critic pair.

\subsection{Novelty and relation to prior models}
We are not aware of a prior model that splits a recurrent vision encoder into priority and value streams on a shared memory, requires that coupling for learning, or requires the grounded stream to be the actor. The closest machine-learning result on actor/critic representation sharing reaches the opposite conclusion for feedforward representations, where separating the actor and critic representations helps. The contrast is informative rather than contradictory: for a recurrent, image-grounded agent whose decision depends on reading a live image at the moment of action, the shared recurrent substrate is what allows credit to reach the actor, and coupling becomes required.

\subsection{Limitations}
The claims are bounded by their scope. The result is a controlled comparison within a single agent family on a single task family, and the brain mapping is an analogy whose direction and substrate differ from the biology as noted above. The collapses are categorical and reproducible across $500$ fresh episodes per configuration, but the present numbers are from trained checkpoints under a stochastic policy rather than an exhaustive seed sweep, and we do not claim that no alternative uncoupled design could ever learn the task --- only that the minimal, parameter-matched uncoupling we tested does not. The functional roles we assign to the two streams are inferred from architecture and behaviour and are made falsifiable below rather than asserted.

\subsection{The network as a generator of falsifiable predictions}
Because the working configuration is interpretable and manipulable, it yields concrete, testable predictions framed as a double dissociation. Clamping or lesioning the value ($q$) circuit should degrade value and context modulation --- the cue-value and reliability effects --- while sparing change detection; lesioning the priority ($\mu$/SC) node should instead collapse detection itself. This maps directly onto SC inactivation abolishing covert detection versus a value-circuit manipulation that leaves detection intact \citep{katz2025optogenetic, krauzlis2018selective}, and it is testable at three levels: behaviourally, the coupling-required and grounded-must-act conditions should hold at training time; representationally, change-evidence should decode preferentially from the priority stream and value/context from the value stream; and causally, perturbing priority attention should shift behaviour and decision threshold while perturbing value attention should shift critic value with behaviour spared. A reward-trained network constrained to be interpretable thus serves not as a claim about the brain but as a generative source of mechanistic hypotheses for systems neuroscience.
